\section{Reglas de prioridad y máxima coincidencia}

\begin{table}[H]
\centering
\small
\caption{Resolución de superposiciones entre patrones.}
\label{tab:prioridad}
\begin{tabularx}{\textwidth}{L{2.7cm}YL{4.3cm}}
\toprule
\tableheader{Caso} & \tableheader{Regla aplicada} & \tableheader{Resultado} \\
\midrule
\texttt{\textbackslash== / \textbackslash= / ==} & Los operadores se prueban por longitud descendente. El lexer no confirma \texttt{\textbackslash=} si existe otro signo igual consecutivo. & \texttt{\textbackslash==} cuando la entrada es \texttt{\textbackslash==}. \\
\rowcolor{palegray} \texttt{=.. / =} & Tras reconocer \texttt{=} como aceptación provisional, se continúa por dos puntos si están disponibles. & \texttt{=..} cuando la entrada es \texttt{=..}. \\
\texttt{:- / :} & Solo \texttt{:-} pertenece al subconjunto. Un \texttt{:} aislado genera un diagnóstico. & \texttt{:-} o error en \texttt{:}. \\
\rowcolor{palegray} \texttt{\_ / \_Variable} & Se reconoce primero el identificador completo. Solo el lexema exacto \texttt{\_} se clasifica como anónimo. & VARIABLE\_ANONIMA o VARIABLE. \\
\texttt{is / island} & Se consume la palabra completa y después se consulta la tabla de operadores alfabéticos. & \makecell[l]{OPERADOR\_\\ARITMETICO o ATOMO.} \\
\rowcolor{palegray} \texttt{-25} & El signo queda fuera del patrón numérico y se emite como operador independiente. & \texttt{-} seguido de ENTERO \texttt{25}. \\
\bottomrule
\end{tabularx}
\end{table}

La implementación aplica estas reglas dentro de un orden de decisión único. Primero reconoce espacios y comentarios, porque no deben producir tokens. Después procesa los literales citados, los números y los identificadores. Los operadores simbólicos se comparan por longitud descendente y los delimitadores se dejan para el final. Si ninguna categoría acepta el carácter actual, se registra un error y se avanza una posición.

Esta secuencia evita dos problemas concretos. Un comentario de bloque debe detectarse antes del operador división, y una palabra completa debe reconocerse antes de decidir si corresponde a \texttt{is}, \texttt{mod} o a un átomo ordinario. La prioridad, por tanto, no reemplaza la máxima coincidencia; ambas reglas trabajan juntas.

\section{Diseño e implementación}

\subsection{Arquitectura de la solución}

\begin{table}[H]
\centering
\small
\caption{Componentes de la implementación.}
\label{tab:arquitectura}
\begin{tabularx}{\textwidth}{L{3cm}L{3.6cm}Y}
\toprule
\tableheader{Archivo} & \tableheader{Componente} & \tableheader{Responsabilidad} \\
\midrule
\path{src/lexer.py} & \texttt{Token} y \texttt{LexError} & Representan salidas inmutables con tipo, lexema, posición y atributo opcional. \\
\rowcolor{palegray} \path{src/lexer.py} & \texttt{SymbolTable} & Interna átomos, variables y literales; reutiliza el índice de un lexema repetido. \\
\path{src/lexer.py} & \texttt{Lexer.scan} & Recorre la fuente una vez, aplica prioridad, actualiza posiciones y recupera errores. \\
\rowcolor{palegray} \path{src/cli.py} & Interfaz CLI & Lee UTF-8 con BOM opcional y emite texto o JSON en UTF-8; retorna 0, 1 o 2 según el resultado. \\
\path{tests/} & Pruebas & Comprueba categorías, posiciones, recuperación, interfaz y corpus completo. \\
\rowcolor{palegray} \path{corpus/} & Entradas completas & Contiene un programa sin errores y otro con varios errores recuperables. \\
\path{scripts/validate.py} & Evidencia reproducible & Ejecuta la suite y el corpus; actualiza resultados JSON y cifras del informe. \\
\bottomrule
\end{tabularx}
\end{table}

La interfaz carga el archivo como texto UTF-8, acepta un BOM inicial opcional y lo entrega a \texttt{Lexer.scan} sin alterar sus saltos de línea. El recorrido mantiene el índice absoluto, la línea y la columna. Antes de consumir un lexema guarda su posición inicial; después actualiza el contador con el texto realmente leído. Los identificadores y números se reconocen mediante patrones compilados, mientras que los literales citados y comentarios de bloque se procesan manualmente para distinguir cierre, salto de línea y escape. Los operadores simbólicos se consultan en una secuencia ordenada desde el lexema más largo. Los códigos de salida son 0 si no hay errores léxicos, 1 si se detectaron errores recuperables o hasta EOF, y 2 ante un archivo ilegible, UTF-8 inválido o argumentos incorrectos.

El resultado agrupa la lista de tokens, la lista de errores y la tabla de lexemas. La interfaz puede imprimir tokens en el formato solicitado o serializar el resultado completo en JSON. Como cada carácter consumido hace avanzar el índice y la cantidad de operadores es fija, el costo del recorrido crece linealmente con el tamaño de la entrada. Los errores consumen el fragmento problemático y el análisis continúa cuando existe un límite de recuperación identificable.

\subsection{Formato de salida}

El siguiente fragmento proviene de \path{corpus/programa_valido.pl}. El quinto campo corresponde al índice de la tabla de lexemas cuando la categoría lo utiliza.

\begin{lstlisting}[style=terminal]
<ATOMO, 'padre', 2, 1, 0>
<PARENTESIS_IZQ, '(', 2, 6>
<ATOMO, 'juan', 2, 7, 1>
<COMA, ',', 2, 11>
<ATOMO, 'ana', 2, 13, 2>
<PARENTESIS_DER, ')', 2, 16>
<PUNTO, '.', 2, 17>
\end{lstlisting}

\subsection{Tabla de lexemas}

La tabla usa diccionarios que conservan un índice por primera aparición. Los átomos no citados y citados comparten una tabla; las variables poseen otra; y los enteros, reales y cadenas se registran como pares \texttt{(tipo, lexema)}, lo que evita mezclar valores escritos de forma distinta. Se conserva la escritura original: \texttt{juan} y \texttt{'juan'} tienen entradas distintas, y no se interpretan los escapes ni se convierten números a valores de Python. Los índices son locales a cada tabla y comienzan en cero; el tipo del token identifica cuál debe consultarse. La variable anónima no se interna porque cada aparición representa una variable nueva y no debe compartir identidad.

\begin{table}[H]
\centering
\small
\caption{Ejemplo de internación sin duplicados.}
\label{tab:lexemas}
\begin{tabularx}{\textwidth}{L{2.5cm}YL{4.5cm}}
\toprule
\tableheader{Tabla} & \tableheader{Ejemplo} & \tableheader{Regla} \\
\midrule
Átomos & \texttt{padre $\rightarrow$ 0; juan $\rightarrow$ 1; ana $\rightarrow$ 2} & El mismo lexema reutiliza el índice. \\
\rowcolor{palegray} Variables & \texttt{X $\rightarrow$ 0; Z $\rightarrow$ 1; Y $\rightarrow$ 2} & \texttt{\_} se excluye. \\
Literales & \texttt{(ENTERO, 52) $\rightarrow$ 0} y \texttt{(CADENA, "Ana Perez") $\rightarrow$ 1} & El tipo y la escritura forman la clave. \\
\bottomrule
\end{tabularx}
\end{table}
